\begin{table}[!htbp]
\caption{Component evidence and attribution limits}\label{tab:6}
\small\renewcommand{\arraystretch}{1.12}\setlength{\tabcolsep}{3pt}
\setlength{\JournalTableWidth}{\dimexpr\linewidth-6\tabcolsep\relax}
\begin{tabular}{@{}>{\raggedright\arraybackslash}p{0.22\JournalTableWidth} >{\raggedright\arraybackslash}p{0.24\JournalTableWidth} >{\raggedright\arraybackslash}p{0.24\JournalTableWidth} >{\raggedright\arraybackslash}p{0.3\JournalTableWidth}@{}}
\toprule
Contrast & Controlled intervention & Observed evidence & Supported interpretation \\
\midrule
S versus G & Add the category prior; keep active geometric planning & Reference: 2 wins / 2 ties; changed background: 1 win / 3 ties & Conditional value of the category channel under the common two-template model \\
X versus Xnf & Enable correction and future diagnosis together & Development mean joint score +2.62\%; extension means positive & Benefit of the correction policy bundle; no isolated posterior or lookahead effect \\
Same-state posterior substitution & Change belief at the same saved decision state & Next action changes in both development configurations & Belief influences the decision; no new endpoint or independent trial \\
Return-feasibility guard & Graph inequality and observed termination records & Qualified runs return; all 32 budget-30 attempts fail the planning gate & Constraint enforcement on this graph; no isolated efficiency gain or physical safety guarantee \\
Shared versus fixed category prior & Change cross-instance sharing in the separate multi-instance study & All nine matched comparisons have identical actions and quality & No measured sharing benefit in the tested conditions \\
\bottomrule
\end{tabular}
\end{table}
